\documentclass[a4paper,10pt]{article}
\include{amsymb}
\include{marvosym}

\usepackage[francais]{babel}
\usepackage[cyr]{aeguill}
\usepackage[applemac]{inputenc}
\usepackage{graphicx}
\usepackage{xspace}
\usepackage[a4paper]{geometry}
\usepackage{latexsym,amsmath,amssymb,textcomp}
\usepackage{moreverb}
\usepackage{listings}
\usepackage{multirow}
\usepackage{titling}
\usepackage{mathabx}

\usepackage{pdfpages}

\usepackage{hyperref}


\newlength{\indentationnota}
\newlength{\largeurlignenota}
\newlength{\paddingnota}
\newlength{\largeurnota}
\setlength{\paddingnota}{5pt}
\setlength{\largeurnota}{0.9cm}

\definecolor{pink}{rgb}{1,0.5,1}
\makeatletter

\newenvironment{pictonote}[1]{% on passe le nom du �chier en argument
\begin{list}{}{%
\setlength{\labelsep}{0pt}%
\setlength{\rightmargin}{15pt}%
\setlength{\paddingnota}{5pt}}
\item%
\setlength{\indentationnota}{%
\@totalleftmargin+\largeurnota+\paddingnota}%
\setlength{\largeurlignenota}{%
\linewidth-\largeurnota-\paddingnota}%
\parshape=3%
\indentationnota\largeurlignenota%
\indentationnota\largeurlignenota%
\@totalleftmargin\linewidth%
\raisebox{-\largeurnota+2.2ex}[0pt][0pt]{%
\makebox[0pt][r]{%
\includegraphics[width=\largeurnota]{#1}%
\hspace{\paddingnota}}}%
\ignorespaces}{%
\end{list}}

\makeatother

\newenvironment{attention}%
{\begin{pictonote}{/Users/benjamin/Documents/Education/LaTeX/danger}}%
{\end{pictonote}}

\newenvironment{unicorn}%
{\begin{pictonote}{/Users/benjamin/Documents/Education/LaTeX/angry_unicorn}}%
{\end{pictonote}}

\newcommand\unicornbox[1]{\colorbox{pink}{\color{white}{#1}}}

\newcommand\cf{\emph{c.f}}

\newcommand\signature{%
\begin{figure}[br]
	\begin{flushright}
	\begin{minipage}{8cm}
		\begin{center}
		\reflectbox{\includegraphics[width=8cm]{/Users/benjamin/Documents/Education/LaTeX/unicorn}}\\
		\theauthor
		\end{center}
	\end{minipage}
	\end{flushright}
\end{figure}}

\newcommand\danger{\raisebox{-0.4ex}{\LARGE $\triangle$ \normalsize} \hspace{-4.1ex}! \hspace{1ex}}
\newcommand\PL{Programmation Logique\xspace}
\newcommand\pl{\bsc{Prolog}\xspace}


\title{RDF - S�ance 2}
\author{Benjamin \bsc{Van Ryseghem}}

\begin{document}
\maketitle
\section{Exercice 1 : R�duction d'une chaine de Freeman}
\subsection{Question 1}

\paragraph{Contour 1}
\verb?"3100077575435443"?

\paragraph{Contour 2}
\verb?"00000767654444213432"?

\subsection{Question 2: R�duction de chaine}

On �limine les couples:
\begin{itemize}
	\item (0,4)
	\item (1,5)
	\item (2,6)
	\item (3,7)
\end{itemize}

\paragraph{Contour 1}
\verb?"55"?

La chaine r�sultante correspond aux transformations reliant l'origine et l'arriv�e.


\paragraph{Contour 2}
\verb?""?

La chaine est vide car le contour est ferm�.

\subsection{Question 3}

Il faut "inverser" la chaine r�sultante.

"55" $\rightarrow$ "11"

\section{Exercice 2 : Transform�e de Hough}
\subsection{Question 1}
\subsection{Question 2}
On d�finit le centre du rep�re en (3.5, 3.5)

On obtient alors: \\
\\
\footnotesize
\begin{tabular}{| c | c c c c c c c c c c c c c c |}

\hline
 & $-3.5\sqrt{2}$& $-3\sqrt{2}$&$-2.5\sqrt{2}$&$-2\sqrt{2}$&$-1.5\sqrt{2}$&$-\sqrt{2}$&$-0.5\sqrt{2}$& $0$ & $0.5\sqrt{2}$ & $\sqrt{2}$ & $1.5\sqrt{2}$ & $2\sqrt{2}$ & $2.5\sqrt{2}$ & $3\sqrt{2}$  \\
$\pi/4$ & 0 & x & x & x & 3 & x & x & x &x&x&x&x&x& x \\
$3\pi/4$ & 0 &x &x&x&x&x&x& x & x & x & x & x & x & x \\
\hline
\end{tabular}
\normalsize


\begin{tabular}{| c | c c c c c c c  c |}
\hline
& -3.5 & -2.5 & -1.5 & -0.5 & 0.5 & 1.5 & 2.5 & 3.5 \\
$0$  x	 & x & x & x & x & x & x & x & x \\
$\pi/2$   x	 & x & x & x & x & x & x & x  & x \\
\hline
\end{tabular}

\signature

\end{document}